\clearpage
\section{Kết luận và hướng phát triển}
\subsection{Kết quả đạt được}
Đồ án đã tích hợp trò chơi suy luận trên lưới với bộ mã hóa CNF, thuật
toán DPLL và Logic Agent. Người chơi và tác tử dùng cùng quy tắc mở thẻ:
kết luận phải suy ra được từ manh mối và trạng thái đã công bố tại lượt
đó. Giao diện hỗ trợ chơi thủ công, gợi ý và giải tự động. Bộ kiểm định
xác nhận cả tính duy nhất của nghiệm khi dùng toàn bộ manh mối lẫn khả
năng mở thẻ tuần tự theo luật chơi.

Trong bộ dữ liệu thực nghiệm đã lưu, cả chín màn đều có nghiệm duy nhất
và Agent hoàn tất cả 45 lượt giải. Thời gian giải trung bình theo kích
thước $3\times3$, $4\times4$ và $5\times5$ lần lượt là 9,12 ms, 28,23 ms
và 115,78 ms. Các phép đo này mô tả bộ màn và cấu hình chạy thực nghiệm.

\subsection{Hạn chế}
Bộ dữ liệu có ba màn cho mỗi kích thước và không ghi nhận lần quay lui
nào trong các lượt giải. Do đó, số liệu thời gian hiện tại chưa thể hiện
chi phí của DPLL trên các màn cần tìm kiếm qua nhiều nhánh. Mã hóa tổ
hợp trực tiếp có thể sinh nhiều mệnh đề khi vùng đếm lớn. Agent cũng
xây lại CNF ở mỗi lần gọi và giải độc lập các truy vấn SAT, kể cả khi
các truy vấn dùng chung phần lớn mệnh đề. Nhật ký lưu các manh mối đang
công khai và kết quả truy vấn, nên người chơi chưa có lời giải thích
chỉ ra tập tiền đề nhỏ đủ để chứng minh một kết luận.

\subsection{Hướng phát triển}
Có thể mở rộng bộ màn bằng các trường hợp buộc DPLL quay lui và các
màn có nghiệm duy nhất nhưng bị kẹt ở trạng thái công khai. Với vùng
đếm lớn, có thể so sánh mã hóa tổ hợp với bộ đếm tuần tự theo số mệnh
đề, số biến phụ và thời gian giải. Một hướng khác là tái sử dụng phần
CNF chung giữa các lượt và dùng giả định tạm trong bộ giải SAT tăng
dần. Khi cải tiến gợi ý, có thể tìm một tập manh mối đủ để chứng minh
nước đi và diễn giải tập đó bằng câu mô tả phù hợp với giao diện.
